\documentstyle[% 


righttag% 


,11pt% 


,amssymb% 


,russian%               remove in Eng. 


,izvru%                 Set izven% in Eng. 


% ,izvru-ks             for brief communications 


]{amsart} 


 


%% ======= For internal remarks 


%% No \No in line 890


%% 


%% 


%% 


 


 


%       Math definitions 


 


\newcommand{\End}{\operatorname{End}} 


\newcommand{\im}{\operatorname{Im}} 


\newcommand{\supp}{\operatorname{supp}} 


\newcommand{\Diff}{\operatorname{Diff}} 


\newcommand{\vraisup}{\operatornamewithlimits{vrai\,sup}} 


\newcommand{\vraiinf}{\operatornamewithlimits{vrai\,inf}} 


% \newcommand{\wh}{\widehat} 


 


\theoremstyle{definition} 


\theorembodyfont{\rm} 


\newtheorem{cornonum}{\hskip\parindent Следствие} 


\renewcommand{\thecornonum}{} 





\renewcommand{\baselinestretch}{1.05}


 


%\hoffset=-15mm                  For Eng. 


 


\setcounter{page}{3} 


\god{1997} 


\nomer{10~(425)} %                Remove in Eng. 


%\nomer{Vol.\,41, No.\,10}        Set in Eng. 


% \str{75-81}                    Set in Eng. 


\udk{517.929} 


\author{\it А.Г.\,Баскаков} 


\title{Спектральный анализ разностных операторов\\ 


с сильно непрерывными коэффициентами}


 


\thanks{Работа выполнена при финансовой поддержке грантов


\No\,NZA000 и 


NZA300 Международного Научного Фонда и Российского фонда 


фундаментальных исследований, проект 95-01-00032.} 


\date{} 


% \pagestyle{myheadings}         Set in Eng. 


 


\pagestyle{plain} %              Remove in Eng. 


\begin{document} 


 


% \thispagestyle{plain}          Set in Eng. 


 


\maketitle 


 


% Body text 


 


Пусть $G$ --- локально компактная абелева группа (с аддитивной формой 


записи 


алгебраической операции), $\wh G$ --- двойственная группа непрерывных 


унитарных 


характеров группы $G$ (с мультипликативной формой записи алгебраической 


операции), $X$ --- комплексное банахово пространство и $\cal F(G,X)$ --- 


банахово пространство измеримых функций, 


определенных на группе $G$ со значениями в $X$ (выбор пространства 


$\cal F(G,X)$ уточняется 


в \S\,1). Через $\End Y$ обозначим банахову алгебру линейных 


ограниченных операторов, действующих в банаховом пространстве~$Y$. 


 


В данной статье изучаются спектральные свойства разностных операторов 


$\cal D\in\End\cal F(G,X)$ 


вида 


\begin{equation} 


(\cal D x)(g)=\sum_{n\ge 1}A_n(g)x(g+g_n),\quad 


x\in\cal F(G,X),\quad g,g_n\in G,\quad n\ge 1, 


\end{equation} 


где функции $A_n$, $n\ge 1$ (называемые коэффициентами оператора $\cal D$), 


принадлежат 


банаховой алгебре $C_s(G,\End X)$ сильно непрерывных и ограниченных на $G$ 


функций 


со значениями в $\End X$. При этом считается выполненным условие 


$\sum\limits_{n\ge 1}\|A_n\|_\infty<\infty$, где 


$\|B\|_\infty=\sup\limits_{g\in G}\|B(g)\|$ $\forall B\in C_s(G,\End X)$. 


Основные результаты статьи связаны с изучением структуры обратных операторов 


к разностным операторам и получены с использованием представлений абелевых 


групп. В теореме 1 доказано, что операторы вида (1) образуют наполненную 


подалгебру в банаховой алгебре $\End\cal F(G,X)$. Она служит основой для 


более глубокого 


исследования спектральных (структурных) свойств разностных операторов. 


Особое внимание уделяется изучению разностных операторов вида 


\begin{equation} 


(\cal D_0x)(g)=x(g)-B(g)x(g-g_0),\quad x\in\cal F(G,X), 


\end{equation} 


где $g_0\in G$ и $B\in C_s(G,\End X)$. Это связано с тем, что абстрактные 


гиперболические операторы 


являются производящими операторами полугрупп операторов вида (2)


(в этом случае $G=\Bbb R$; см. теоремы 


4--6). Теорема 6 позволяет во многих вопросах свести изучение абстрактных 


гиперболических операторов к изучению соответствующих разностных 


операторов. По всей видимости, этот результат может оказаться полезным при 


исследовании задач управления. Спектральные свойства операторов вида (2) 


рассматриваются в теоремах 2 и 3. 


 


\section{Наполненность подалгебры разностных операторов} 


 


Символом $\cal F(G,X)$ обозначим одно из следующих банаховых пространств 


измеримых 


(по Бохнеру) функций, определенных на группе $G$ со значениями в банаховом 


пространстве $X$; $L_p=L_p(G,X)$ --- пространство суммируемых со степенью 


$p\in[1,\infty)$ (относительно 


меры Хаара $\mu$ на $G$) функций, $L_\infty=L_\infty(G,X)$ --- пространство 


существенно ограниченных 


на $G$ функций


($\|x\|_{\infty}=\vraisup\limits_{g\in G}\|x(g)\|$, $x\in L_\infty$), 


$C=C(G,X)$ --- подпространство непрерывных функций из $L_\infty(G,X)$, 


$C_0=C_0(G,X)$ --- 


подпространство функций из $C(G,X)$, сходящихся к нулю на бесконечности. 


Если $X=\Bbb C$ 


--- поле комплексных чисел, то введенные банаховы пространства обозначаются 


символами $L_p(G)$, $p\in[1,\infty)$, $C(G)$, $C_0(G)$. При этом $L_1(G)$ 


является 


коммутативной банаховой алгеброй со сверткой функций в качестве умножения. 


 


Вместо символа $\cal F(G,X)$ часто будет использоваться символ $\cal F$. То 


же замечание относится и к конкретным функциональным пространствам. 


 


В статье систематически используются следующие два изометрические 


представления $S:G\to\End\cal F$, $V:\wh G\to\End\cal F$, определенные 


равенствами 


$$ 


(S(g)x)(u)=x(u+g),\quad (V(\gamma)x)(u)=\gamma(u)x(u),\quad 


u,g\in G,\quad \gamma\in\wh G,\quad x\in\cal F. 


$$ 


Используемые далее понятия и результаты из теории векторных почти 


периодических функций можно найти в монографиях \cite{Lev}--\cite{Dun}.


 


\begin{defn} Линейный оператор $A\in\End\cal F$ назовем разностным, если 


функция $\wh A:\wh G\to\End\cal F$, определенная 


равенствами $\wh A(\gamma)=V(\gamma^{-1})AV(\gamma)$, $\gamma\in\wh G$, 


является непрерывной почти периодической (п.\,п.) функцией. 


\end{defn} 


 


Если оператор $\cal D\in\End\cal F$ определен формулой (1), то функция 


$\wh D:\wh G\to\End\cal F$ 


имеет вид 


$\wh{\cal D}(\gamma)=V(\gamma^{-1})\cal D V(\gamma)=


\sum\limits_{n\ge 1}\widetilde A_nS(g_n)\gamma(g_n)$, где 


$\widetilde A_n$ --- 


оператор умножения на функцию $A_n$. Из равенств 


$\|\widetilde A_n\|=\|A_n\|_\infty$, $n\ge 1$, 


следует, что $\wh{\cal D}$ --- п.\,п. 


функция. Таким образом, $\cal D$ --- разностный оператор (в смысле 


определения 1). 


 


В частности, если $G\in\Bbb Z$ --- группа целых чисел и 


$l_p(\Bbb Z)=L_p(\Bbb Z)$, $p\in[1,\infty)$, то каждый 


оператор $A\in\End l_p(\Bbb Z)$, 


определяемый матрицей $(a_{ij};\;i,j\in\Bbb Z)$ со свойством 


$\sum\limits_{n\ge 1}\sup\limits_{i-j=n}|a_{ij}|<\infty$, является 


разностным. Он представим 


в виде (1), где $g_n=n$, $n\in\Bbb Z$, и функции 


$A_n:\Bbb Z\to\Bbb C$, $n\in\Bbb Z$, определяются 


равенствами $A_n(k)=a_{k(k+n)}$, $k\in\Bbb Z$. Нетрудно 


доказать, что оператор $\cal D\in\End l_p(\Bbb Z)$ является разностным 


тогда и только тогда, 


когда он является пределом в $\End l_p(\Bbb Z)$ операторов с конечным 


числом диагоналей. 


Следовательно, $\Diff l_p(\Bbb Z)$ содержит наиболее важные в приложениях 


операторы из $\End l_p(\Bbb Z)$. 


 


\begin{defn} Пусть $A\in\Diff\cal F$ и п.\,п. функция 


$\wh A:\wh G\to\End\cal F$ имеет ряд Фурье вида 


$\widetilde A(\gamma)\sim\sum\limits_{n\ge 1}A_n\gamma(g_n)$. Ряд 


$\sum\limits_{n\ge 1}A_n$ назовем рядом 


Фурье разностного оператора $A$, а множество 


$\{g_1,g_2,\dotsc\}\subset G$ показателей Фурье функции $\wh A$ 


обозначим символом $\Lambda(A)$. 


\end{defn} 


 


Из определения коэффициентов Фурье п.\,п. функции следует, что операторы 


$A_n\in\End\cal F$, $n\ge 1$,


могут быть вычислены по формуле 


\begin{equation} 


A_n=\int_{\overline G}\overline A(\tau)\tau(-g_n)\overline\mu(d\tau), 


\end{equation} 


где $\overline G$ --- боровская компактификация группы $\wh G$ 


(т.\,е. $\overline G=\wh G_d$ --- 


двойственная 


к группе $G_d$, $G_d$ --- группа $G$, наделенная дискретной топологией), 


$\overline\mu$ --- мера 


Хаара на компактной группе $\overline G$ и 


$\overline A:\overline G\to\End\cal F$ --- непрерывное продолжение 


функции $\wh A$ на $\overline G$. Отметим свойство 


$\lim\limits_{n\to\infty}\|A_n\|=0$. 


 


Символом $\Diff_V\cal F$ или $\End_V\cal F$ обозначим множество операторов 


из алгебры 


$\End\cal F$, перестановочных с операторами представления 


$V:\wh G\to\End\cal F$. Ясно, что 


$\Diff_V\cal F\subset\Diff\cal F$ и $\Diff_V\cal F$ --- 


замкнутая подалгебра из алгебры $\End\cal F$. Алгебра $\Diff\cal F$ 


содержит операторы умножения на 


функции из алгебры $C_s(G,\End X)$. Однако не всякий оператор из 


$\Diff_V\cal F$ является 


оператором умножения (см. \cite{Kur}, гл.\,II). 


 


Подалгебра $\cal A$ из банаховой алгебры $\cal B$ называется наполненной 


\cite{Bou}, если каждый 


обратимый в $\cal B$ элемент из алгебры $\cal A$ обратим и в $\cal A$. 


 


\begin{lem} Множество разностных операторов $\Diff\cal F$ образует 


наполненную замкнутую подалгебру 


в банаховой алгебре $\End\cal F$. Если 


$A\sim\sum\limits_{n\ge 1}A_n$ --- ряд Фурье оператора 


$A\in\Diff\cal F$, то 


операторы 


$A_n$, $n\ge 1$, являются разностными, $\Lambda(A_n)=\{g_n\}$ --- 


одноточечное множество и они допускают представление вида 


\begin{equation} 


A_n=A_n^0S(g_n),\quad A^0_n\in\Diff_V\cal F. 


\end{equation} 


\end{lem} 


 


\begin{pf} Непосредственно из свойств п.\,п. функций следует, что 


$\Diff\cal F$ --- 


замкнутая 


наполненная подалгебра из $\End\cal F$. Например, если 


$A\in\Diff\cal F$ --- обратимый в 


алгебре $\End\cal F$ 


оператор и $B=A^{-1}$, то функция $\wh B:\wh G\to\End\cal F$ может быть 


представлена 


в виде $\wh B(\gamma)=(\wh A(\gamma))^{-1}$, $\gamma\in\wh G$, и поэтому 


является почти периодической. Таким образом, 


$A^{-1}\in\Diff\cal F$. 


 


Непосредственно из формулы (3) следует, что имеют место равенства 


$V(\gamma^{-1})A_nV(\gamma)=\gamma(g_n)A_n$, $n\ge 1$. 


Поэтому $A_n\in\Diff\cal F$ и $\Lambda(A_n)=\{g_n\}$ $\forall n\ge 1$. 


 


Положим $A_n^0=A_nS(-g_n)$, $n\ge 1$. Из равенств 


$V(\gamma^{-1})S(g_n)V(\gamma)=\gamma(-g_n)S(-g_n)$, $\gamma\in\wh G$, 


следует, что 


$V(\gamma^{-1})A^0_nV(\gamma)=A^0_n$ $\forall\gamma\in\wh G$, 


$\forall n\ge 1$. 


Следовательно, $A^0_n\in\Diff_V\cal F$. 


\end{pf} 


 


\begin{cornonum} $\Diff_V\cal F$ --- замкнутая наполненная подалгебра из 


$\End\cal F$. 


\end{cornonum} 


 


Операторы $A^0_n\in\Diff_V\cal F$, $n\ge 1$, из представления (4) назовем 


коэффициентами разностного оператора $A$. 


 


Ряд последующих утверждений (леммы 2--6) будет тесно связан со структурой 


операторов из алгебры $\Diff_V\cal F$ (коэффициентами разностных 


операторов). 


 


Далее символом $C_{0,k}(G)$ обозначена подалгебра функций из алгебры 


$C_0(G)=C_0(G,\Bbb C)$, 


имеющих 


компактный носитель. Через $V(f)$ обозначим оператор из $\Diff\cal F$ 


умножения на функцию 


$f\in C(G)=C(G,\Bbb C)$. Последовательность функций $(f_n)$ из алгебры 


$C_{0,k}(G)$ называется ограниченной 


аппроксимативной единицей (о.\,а.\,е.), если выполнены условия:


1)~$\lim\limits_{n\to\infty}f_n(0)=1$, 


2)~$\lim\limits_{n\to\infty}\|S(g)f_n-f_n\|_\infty=0$ $\forall g\in G$. 


 


\begin{defn} Линейный оператор $A\in\End\cal F$ назовем $c$-непрерывным 


(ср.\cite{Muh} и \cite{Slu}), если для любой 


о.\,а.\,е. $(f_n)$ выполнено соотношение 


$\lim\limits_{n\to\infty}\|V(f_n)A-AV(f_n)\|=0$. 


\end{defn} 


 


Множество всех $c$-непрерывных операторов из алгебры $\End\cal F$ обозначим 


символом $\End_c\cal F$. 


Ясно, что $\cal D\in\End_c\cal F$, если $\cal D$ определен формулой (1). 


Непосредственно из определения 


3 следует, что $\End_c\cal F$ и 


$\Diff_c\cal F=\End_c\cal F\cap\Diff\cal F$ --- замкнутые подалгебры из 


$\End\cal F$ и $\Diff\cal F$ соответственно. 


 


\begin{lem} $\End_c\cal F$ и $\Diff_c\cal F$ --- наполненные подалгебры из 


$\End\cal F$. 


\end{lem} 


 


\begin{pf} Если $A\in\End_c\cal F$ обратим в $\End\cal F$, то 


$$ 


\lim_{n\to\infty}\|V(f_n)A^{-1}-A^{-1}V(f_n)\|= 


\lim_{n\to\infty}\|A^{-1}(AV(f_n)-V(f_n)A)A^{-1}\|=0 


$$ 


для любой о.\,а.\,е. $(f_n)$. Если $A\in\Diff_c\cal F$, то следует 


использовать лемму 1. 


\end{pf} 


 


Через $\cal F_0=\cal F_0(G,X)$ обозначим замыкание в $\cal F(G,X)$ функций 


из $\cal F(G,X)$, имеющих компактный 


носитель. Ясно, что $\cal F_0=\cal F$, если 


$\cal F=L_p$, $p\in[1,\infty)$, или $\cal F=C_0$. Отметим, что 


$\lim\limits_{n\to\infty}V(f_n)x=x$ $\forall x\in\cal F_0$,


$\forall\text{ о.\,а.\,е. } (f)$. 


 


\begin{lem} Если $A\in\End_V\cal F\cap\End_c\cal F$, то $\cal F_0$ --- 


инвариантное для $A$ подпространство, и имеют место равенства 


\begin{equation} 


AV(\varphi)=V(\varphi)A,\quad \varphi\in C_0(G). 


\end{equation} 


\end{lem} 


 


{\bf Доказательство.} Если $x_0\in\cal F_0$, то 


$Ax_0=\lim\limits_{n\to\infty}AV(f_n)x_0=V(f_n)Ax_0\in\cal F_0$. 


 


Пусть функция $\varphi\in C_{0,k}(G)$ имеет преобразование Фурье 


$\wh\varphi$, 


принадлежащее алгебре 


$L_1(\wh G)$. Если $x_0\in\cal F_0$, то функции 


$\gamma\mapsto V(\gamma)x_0$, 


$\gamma\mapsto V(\gamma)Ax_0:\wh G\to\cal F$ непрерывны. 


Интегрируя равенства 


$\wh\varphi(\gamma)V(\gamma)Ax_0=\wh\varphi(\gamma)AV(\gamma)x_0$, 


$\gamma\in\wh G$, 


получаем 


$$ 


V(\varphi)Ax_0=\int_{\wh G}\wh\varphi(\gamma)V(\gamma)Ax_0d\gamma= 


\int_{\wh G}A\wh\varphi(\gamma)V(\gamma)x_0d\gamma=AV(\varphi)x_0. 


$$ 


В силу плотности таких функций $\varphi$ в $C_0(G)$ равенства (5) имеют 


место на $\cal F_0$. Если $x\in\cal F$, то для любой о.\,а.\,е. $(f_n)$


выполняются равенства 


$$ 


V(\varphi)Ax-AV(\varphi)x = \lim_{n\to\infty}(V(f_n)V(\varphi)Ax- 


V(f_n)AV(\varphi)x)= 


V(\varphi)\lim_{n\to\infty}(V(f_n)Ax-AV(f_n)x)=0.\quad\square 


$$ 


 


\begin{cor} Пусть $A\in\End_V\cal F\cap\End_c\cal F$. Тогда 


\begin{itemize} 


\item[1)] $\supp Ax\subset\supp x$ $\forall x\in\cal F$; 


\item[2)] $Ax=Ay$ на открытом множестве $U\subset G$, если $x=y$ на $U$; 


\item[3)] $\|(Ax)(g)\|\le\|A\|\,\|x(g)\|$ $\forall x\in C_0$, 


$\forall g\in G$, если $\cal F=L_\infty$ и $AC_0\subset C_0$; 


\item[4)] $A\cal F_0\subset\cal F_0$, если $\cal F=L_\infty$. 


\end{itemize} 


\end{cor} 


 


\begin{cor} Если $A,B\in\End_V\cal F\cap\End_c\cal F$ совпадают на 


$\cal F_0$, то $A=B$. 


\end{cor} 


 


\begin{lem} Если $A\in\End L_p$, $p\in[1,\infty)$, то 


$Ax\in L_\infty$ и 


$\|Ax\|_\infty\le\|A\|\,\|x\|_\infty$ $\forall x\in L_\infty$ с компактным 


носителем $\supp x$. 


\end{lem} 


 


\begin{pf} Пусть $x\in L_\infty$, $\|x\|_\infty\le 1$ и $\supp x$ --- 


компакт. Предположим, что $\|A\|<\|Ax\|_\infty=\alpha\le\infty$. Тогда 


существуют компакт $K\subset G$ и 


открытое множество $U\subset G$ такие, что 


1)~$K\subset G$, 


2)~$\vraiinf\limits_{g\in K}\|y(g)\|>\|A\|$, где $y=Ax$, 


3)~$\mu(K)>\mu(U)\|A\|^p\alpha^{-p}$. Выберем 


функцию $\varphi\in C_0(G)$ такую, чтобы 


$0\le\varphi\le1$, $\varphi\equiv 1$ на $K$ и $\supp\varphi\subset U$. 


Тогда $\|\varphi x\|_p\le\mu(U)^{1/p}$ ($\|\cdot\|_p$ --- норма в $L_p$). 


 


Функция $\varphi y=\varphi Ax=A(\varphi x)$ допускает следующую оценку 


нормы $\|\varphi y\|_p\le\|A\|\mu(U)^{1/p}$. С другой стороны, 


$\|\varphi y\|_p>\alpha\mu(K)^{1/p}>\|A\|\mu(U)^{1/p}$. Получено 


противоречие. 


\end{pf} 


 


\begin{lem} Любой оператор $A\in\End_V\cal F\cap\End_c\cal F$ допускает 


единственное расширение $\overline A$ на $L_\infty$, если 


$\cal F=L_p$, $p\in[1,\infty)$ 


(и на $C$, если $\cal F=C_0$), причем 


$\|\overline A\|\le\|A\|$ и 


$\overline A\in\End_V L_\infty\cap\End_c L_\infty$ 


{\em (}$\overline A\in\End_V C\cap\End_c C$, если $\cal F=C_0${\em )}. 


\end{lem} 


 


\begin{pf} Пусть $\cal F=L_p$ (случай $\cal F=C_0$ рассматривается 


аналогично). Функция $\overline Ax$, 


где $x\in L_\infty$, определяется на любом открытом множестве с компактным 


замыканием $\overline U$ 


следующим образом: 


$(\overline Ax)(g)=(A\varphi x)(g)$, где $\varphi$ --- любая функция из 


$C_0(G)$ с компактным носителем 


$\supp\varphi$ и $\varphi=1$ на $\overline U$. Непосредственно из леммы 3, 


ее следствий и леммы 4 следует 


корректность определения оператора $\overline A\in\End L_\infty$ и оценка 


$\|\overline A\|\le\|A\|$. Поскольку 


имеют место равенства 


$V(\gamma^{-1})\overline AV(\gamma)=\overline A$, 


$\gamma\in\wh G$, $\overline{V(f)A}=V(f)\overline A$, 


$\overline{AV(f)}=\overline AV(f)$, 


$f\in C_{0,k}(G)$, 


то для любой о.\,а.\,е. $(f_n)$ получаем 


$\|V(f_n)\overline A-\overline AV(f_n)\|\le 


\|V(f_n)A-AV(f_n)\|\to 0$ при $n\to\infty$. 


\end{pf} 


 


\begin{lem} Если $\cal F\ne L_\infty$ и для оператора 


$A\in\End_V\cal F\cap\End_c\cal F$ выполнено условие 


\begin{equation} 


\lim_{h\to\infty}\|S(h)AS(-h)x-Ax\|_\infty=0 


\quad \forall x\in C_0(G,X), 


\end{equation} 


то он представим в виде $(Ax)(g)=A_0(g)x(g)$, $x\in\cal F$, где 


$A_0\in C_S(G,\End X)$ . 


\end{lem} 


 


\begin{pf} Поскольку 


$S(h)V(\gamma)=\gamma(h)V(\gamma)S(h)$ $\forall h\in G$, 


$\forall\gamma\in\wh G$, то 


$S(h)AS(-h)\in\End_V\cal F\cap\End_c\cal F$ 


и в силу 


леммы 4 


$S(h)AS(-h)Ax-Ax\in L_\infty$ $\forall x\in C_0(G,X)$. Из леммы 5 получаем 


существование оператора $\overline A\in\End L_\infty$ 


(если $\cal F=L_p$, $p\in[1,\infty$)\,), 


принадлежащего алгебре 


$\End_V L_\infty\cap\End_c L_\infty$. Докажем, что 


$\overline AC_0\subset C_0$ при $\cal F=L_p$ (случай $\cal F=C$ 


очевиден). 


 


Для этого рассмотрим направленность $\frak A=\{U\}$ компактных окрестностей 


нуля группы $G$ 


(направленных по включению множеств) и направленность операторов 


$\{\overline A_U,\;U\in\frak A\}$ из 


алгебры $\End C_0$, имеющих вид 


\begin{equation} 


\overline A_U x=\frac{1}{\mu(U)}\int_U S(h)\overline AS(-h)x\,dh, 


\quad x\in C_0. 


\end{equation} 


Докажем, что $\overline A_Ux\in C_0$ $\forall x\in C_0$. Пусть вначале 


функция $x\in C_0(G,X)$ имеет вид $x(g)=\varphi(g)x_0$, где 


$\varphi\in C_0(G)$, $x_0\in X$. Тогда из 


(7) получаем 


$\overline A_Ux=\varphi_0\mu(U)^{-1}\int\limits_US(h) 


\overline A\overline x_0dh\in C_0(G,X)$, где 


$\overline x_0(g)=x_0$ $\forall g\in G$. В силу 


плотности таких функций $x$ в $C_0(G,X)$ получаем включение 


$\overline A_UC_0\subset C_0$. 


 


Из условия (6) следует, что для любой функции $x\in C_0$ 


$$ 


\lim_U\|\overline Ax-\overline A_Ux\|_\infty \le 


\lim_U\|\mu(U)^{-1}\|\int_U\|S(h)\overline AS(-h)x-Ax\|_\infty 


dh=0. 


$$ 


Поэтому $\overline Ax\in C_0$ $\forall x\in C_0$ и, следовательно, 


$\overline AC\subset C$, причем 


$\|\overline A\|_{\End C}\le \|A\|$ (см. леммы 4 и 5). 


 


Введем в рассмотрение функцию $A_0\in C_S(G,\End X)$, положив 


$A_0(g)x_0=(\overline A\overline x_0)(g)$, где 


$\overline x_0(g)\equiv x_0\in X$, и 


пусть $B\in\End\cal F$ --- оператор 


умножения на $A_0$. Непосредственно из определения операторов 


$\overline A$, $B$ и лемм 3, 


5 получаем, что 


$(\overline A(\varphi_0x_0))(g)=\varphi_0(g)\overline A 


(\overline x_0)=\varphi_0(g)A_0(g)x_0=(B(\varphi_0x_0))(g)$ 


$\forall g\in G$. Таким 


образом, $\overline A=B$ на $C_0$. Поскольку $C_0$ плотно в $L_p$ 


($p\in[1,\infty)$), 


то $A=B$ при $\cal F=L_p$. При $\cal F=C$ следует использовать лемму 5. 


\end{pf} 


 


Пусть $\sigma$ --- счетная подгруппа из группы $G$ и 


$\alpha:\sigma\to\Bbb R_+=\{t\in\Bbb R:t\ge 0\}$ --- весовая функция 


(т.\,е. 


$\alpha(g)\ge 1$, $\alpha(g_1+g_2)\le\alpha(g_1)\alpha(g_2)$ 


$\forall g,g_1,g_2\in G$), 


удовлетворяющая условию 


$\lim\limits_{n\to\infty}n^{-1}\ln\alpha(ng)=0$ $\forall g\in\sigma$. 


Символом 


$\Diff(\cal F,\alpha)$ обозначим подалгебру разностных операторов вида (1), 


для которых $\Lambda(\cal D)=\{g_1,g_2,\dotsc\}\subset\sigma$ и 


$\sum\limits_{n\ge 1}\|A_n\|_\infty\alpha(g_n)<\infty$. 


 


\begin{thm} Пусть $\cal D\in\Diff(\cal F,\alpha)$ --- обратимый оператор. 


Тогда $\cal D^{-1}\in\Diff(\cal F,\alpha)$, т.\,е. он имеет вид 


$$ 


(\cal D^{-1}x)(g)=\sum_{k\ge 1}B_k(G)X(g+h_k), 


$$ 


где $B_k\in C_S(G,\End X)$, $k\ge 1$, $h_k\in\sigma$ $\forall k\ge 1$ 


и $\sum\limits_{k\ge 1}\|B_k\|\alpha(h_k)<\infty$. 


\end{thm} 


 


\begin{pf} Функция $\wh{\cal D}(\gamma)=V(\gamma^{-1})\cal DV(\gamma)$, 


$\gamma\in\wh G$ ($\wh{\cal D}:\wh G\to\End\cal F$) --- п.\,п.


функция и поэтому будет п.\,п. функцией 


$\wh{\cal D^{-1}}(\gamma)=V(\gamma^{-1})\cal D^{-1}V(\gamma)$, 


$\gamma\in\wh G$, т.\,е. 


$\cal D^{-1}$ --- разностный оператор. Из


(\cite{Bask1}, лемма 1) следует, что оператор 


$\cal D^{-1}$ допускает представление


\begin{equation} 


\cal D^{-1}=\sum_{k\ge 1}\widetilde B_kS(h_k),\quad h_k\le\sigma, 


\end{equation} 


где $\widetilde B_k\in\End_V\cal F$ (см. лемму 1), 


$\sum\limits_{k\ge 1}\|B_k\|\alpha(h_k)<\infty$. Поскольку 


$\wh{\cal D^{-1}}:G\to\End\cal F$ --- 


непрерывная п.\,п. функция, 


значениями которой в силу леммы 2 являются операторы из алгебры 


$\End_c\cal F$, то из 


формулы (3) для коэффициентов Фурье 


$\widetilde B_kS(h_k)$, $k\ge 1$, оператора $\cal D^{-1}$ следует, что 


$\widetilde B_kS(h_k)\in\Diff_C\cal F$ $\forall k\ge 1$. Следовательно, 


$\widetilde B_k\in\End_V\cal F\cap\End_c\cal F$, $k\ge 1$. 


 


Вначале предположим, что $\cal F=L_p$, $p\in[1,\infty)$. Поскольку 


$\|\widetilde B_kS(h_k)\|=\|\widetilde B_k\|$, то из 


леммы 5 следует, что 


$\widetilde B_k$ допускает расширение на $L_\infty$ с нормой, не 


превосходящей $\|\widetilde B_k\|$. Отсюда 


в силу абсолютной сходимости ряда Фурье для $\cal D^{-1}$ получаем, что 


$\cal D^{-1}$ определен 


и ограничен в $L_\infty$, причем из следствия леммы 3 и представления (8) 


следует включение $\cal D^{-1}((L_\infty)_0)\subset(L_\infty)_0$. Поэтому 


для любой функции $x\in C_{0,k}(G,X)$ получаем при $h\to 0$ 


\begin{gather*} 


\|(S(h)\cal D^{-1}S(-h)-\cal D^{-1})x\|_\infty= 


\|S(h)\cal D^{-1}S(-h)(\cal Dx-S(h)\cal DS(-h))\cal D^{-1}x\|_\infty\le\\ 


% 


\le\|\cal D^{-1}\|_{\End L_\infty}\|(\cal Dx-S(h)\cal DS(-h)) 


\cal D^{-1}x\|_\infty\to 0, 


\end{gather*} 


если учесть, что $\cal D^{-1}x\in(L_\infty)_0$. Такое же соотношение имеет 


место при $\cal F=C$. 


 


Еще раз используя представление (3), получаем, что для операторов 


$\widetilde B_k$, $k\ge 1$, выполнено условие (6) леммы 6 и, следовательно, 


для них выполнены все 


условия этой леммы. Поэтому операторы $\widetilde B_k$, $k\ge 1$, 


являются операторами умножения на некоторые функции 


$B_k\in C_s(G,\End X)$ при $\cal F=L_p$, $p\in[1,\infty)$, 


$\cal F=C,C_0$. 


 


Пусть $\cal F=L_\infty$. Поскольку $\widetilde B_kC_0\subset C_0$ (см. 


доказательство леммы 6), то $\cal D^{-1}C_0\subset C_0$ и поэтому сужение 


$\cal D_0$ оператора $\cal D$ на $C_0$ является обратимым оператором. 


Значит, $\widetilde B_k\in\End C_0$, $k\ge 1$, являются 


операторами умножения на некоторые функции $B_k\in C_S(G,\End X)$. В силу 


плотности $C_0$ в $L_1$ и из 


вида операторов $\cal D_0$ и $\cal D_0^{-1}$ следует, что они допускают 


ограниченное расширение на 


$L_1$, причем подпространство $L_1\cap L_\infty$ переводят в себя. Поэтому 


$\cal D_0^{-1}=\cal D^{-1}$ и 


$(\widetilde B_kx)(g)=B_k(g)x(g)$, $k\ge 1$, на $L_1\cap L_\infty$. В 


силу следствия 2 леммы 3 это представление операторов $\widetilde B_k$ 


имеет место и на $L_\infty$. 


\end{pf} 


 


\begin{cornonum} Спектр разностного оператора $\cal D$ вида (1) не зависит 


от выбора пространства $\cal F(G,X)$. 


\end{cornonum} 


 


Отметим, что утверждение теоремы 1 в случае $\alpha(g)\equiv 1$, 


$X$ --- 


конечномерное 


пространство и $A_n\in L_\infty(G,X)$, $n\ge 1$, приведено в работе 


(\cite{Kur}, гл.\,II). Там же излагается история вопроса. 


 


\section{Спектральные свойства оператора взвешенного сдвига} 


 


Пусть $g_0$ --- такой элемент из группы $G$, что множество 


$\{\gamma(g_0),\;\gamma\in\wh G\}$ плотно на 


окружности $\Bbb T=\{\lambda\in\Bbb C:|z|=1\}$. 


Рассмотрим разностный оператор (оператор взвешенного сдвига) 


$A\in\End\cal F(G,x)$ вида 


$$ 


(Ax)(g)=B(g)x(g-g_0),\quad x\in\cal F=\cal F(G,X), 


$$ 


где $B\in C_S(G,\End X)$ (подробные комментарии к истории исследования 


операторов взвешенного сдвига имеются в монографии \cite{Ant}). 


 


Предположим, что число 1 не принадлежит спектру $\sigma(A)$ оператора $A$. 


Следующие равенства 


\begin{equation} 


V(\gamma)AV(\gamma^{-1})=\gamma(g_0)A,\quad \gamma\in\wh G,


\end{equation} 


означают подобие оператора $A$ операторам вида $\gamma(g_0)A$, 


$\gamma\in\wh G$. Поскольку 


подобные 


операторы имеют одинаковые спектры, то из открытости резольвентного 


множества $\rho(A)=\Bbb C\setminus\alpha(A)$ и равенств (9) следует, что 


выполнено условие 


\begin{equation} 


\sigma(A)\cap\Bbb T=\emptyset. 


\end{equation} 


 


\begin{thm} Пусть разностный оператор $\cal D=I-A\in\End\cal F$ обратим. 


Тогда имеют место следующие свойства{\em :} 


\begin{itemize} 


\item[1)] $\sigma(A)=\sigma_-\cup\sigma_+,$ 


$\sigma_-=\{\lambda\in\sigma(A):|\lambda|<1\},$ 


$\sigma_+=\{\lambda\in\sigma(A):|\lambda|>1\};$ 


\item[2)] проектор Рисса $P_-=P(\sigma_-,A)$, построенный по спектральному 


множеству $\sigma_-$, является оператором умножения на проекторозначную 


функцию $P\in C_S(G,\End X);$ 


\item[3)] оператор $\cal D^{-1}$ представим в виде 


$\cal D^{-1}=\sum\limits_{k\in\Bbb Z}B_k$, где 


$(B_kx)(g)=\cal D_k(g)x(g-kg_0)$, 


$\cal D_k\in C_S(G,\End X)$, 


причем 


операторы $B_k\in\End\cal F$, $k\in\Bbb Z$, являются 


голоморфными функциями от оператора $A$ и определяются из соотношений 


\begin{equation} 


B_k=A^{-k}P_-,\quad k\le 0;\quad B_kA^k=A^kB_k=-P_+;\quad 


B_kP_-=P_-B_k=0,\quad k>0, 


\end{equation} 


где $P_+=I-P_-$ --- дополнительный к $P_-$ проектор{\em ;} 


\item[4)] для любых чисел $\delta_-$, $\delta_+$ из интервала $(0,1)$, 


удовлетворяющих условиям $\delta_\pm\|\cal D^{-1}\|<1$, справедливы оценки 


\begin{gather} 


\|P\|_\infty\le \|\cal D^{-1}\|,\quad \|B_k\|=\|\cal D_k\|_\infty \le 


\|\cal D^{-1}\|,\quad k\in\Bbb Z,\\ 


% 


\|B_k\|\le\|\cal D^{-1}\|(1-\delta_-\|\cal D^{-1}\|)^{-1} 


(1-\delta_-)^{-k}, \quad k\le -1,\\ 


% 


\|B_k\|\le\|\cal D^{-1}\|(1-\delta_+\|\cal D^{-1}\|)^{-1} 


(1+\delta_+)^{-k}, \quad k\ge -1,


\end{gather} 


$$ 


\sum_{k\in\Bbb Z}\|B_k\|\le \|\cal D^{-1}\|(1+8\|\cal D^{-1}\|). 


$$ 


\end{itemize} 


\end{thm} 


 


\begin{pf} Свойство 1) следует из условия (10). Имеет место равенство 


$\cal F=\cal F_-\oplus\cal F_+$, где 


$\cal F_\pm=\im P_\pm$ 


--- образы проекторов $P_\pm$, и оператор $A$ имеет вид 


$A=A_-\oplus A_+$, где $A_\pm=A\mid\cal F_\pm$ --- сужение 


$A$ на 


$\cal F_\pm$. Тогда $\sigma(A_\pm)=\sigma_\pm$ и 


$\cal D^{-1}=(I_--A_-)^{-1}\oplus(I_+-A_+)^{-1}$, $I_\pm$ --- тождественный 


оператор в $\cal F_\pm$. Поскольку спектральные 


радиусы $r(A_-)$ и $r(A_+^{-1})$ операторов $A_-$ и $A_+^{-1}$ меньше 1, то 


имеет место представление 


\begin{equation} 


\cal D^{-1}=\biggl(\,\sum_{k\ge 0}A_-^k\biggr)\oplus 


\biggl(\,\sum_{k\le -1}A_+^k\biggr)=\sum_{k\in\Bbb Z}B_k, 


\end{equation} 


где 


$B_k=A_-^k\oplus 0=A^kP_-$, $k\ge 0$; $B_k=0\oplus A_+^k$, $k\le -1$. 


Используя простейшие свойства функционального исчисления 


(см. \cite{Dun}, гл.\,VII), операторы $B_k$, $k\in\Bbb Z$, можно 


представить формулами 


\begin{equation} 


B_k=\frac{1}{2\pi i}\int_{\Bbb T}\lambda^{-k}R(\lambda,A)d\lambda, 


\quad k\in\Bbb Z, 


\end{equation} 


из которых следуют соотношения (11). Из (16) получаются равенства 


\begin{gather*} 


\wh B_k(\gamma)=V(\gamma^{-1})B_kV(\gamma)=\frac{1}{2\pi i} 


\int_{\Bbb T}\lambda^{-k}R(\lambda,\gamma(g_0)A)d\lambda=\\ 


=\frac{1}{2\pi i}\int_{\Bbb T}\lambda^{-k}\gamma(-g_0) 


R(\lambda\gamma(-g_0),A)d\lambda=\gamma(-kg_0)B_k. 


\end{gather*} 


Таким образом, $\wh B_k$, $k\in\Bbb Z$, --- п.\,п. функции и потому они 


являются разностными 


операторами с $\Lambda(B_k)=-\{kg_0\}$. Из теоремы 1 получаем (см. ее 


доказательство), 


что операторы $B_k$, $k\in\Bbb Z$, допускают представление вида 


$(B_kx)(g)=\cal D_k(g)x(g-kg_0)$, $k\in\Bbb Z$, где 


$\cal D_k\in C_S(G,\End X)$. 


Поскольку $B_0=P_-$, то 


$(P_-x)(g)=P(g)x(g)$, $x\in\cal F$, 


где $P(g)=\cal D_0(g)$, $g\in G$, --- проекторозначная функция из 


$C_S(G,\End X)$. Итак, доказаны свойства 


1)--3). Оценки (12) следуют из оценок коэффициентов Фурье п.\,п. функции 


($\|B_k\|\le\sup\limits_{\gamma\in\wh G} 


\|\wh{\cal D^{-1}}(\gamma)\|=\|\cal D^{-1}\|$). 


 


Получим оценки (13). Если $\sigma_-=\emptyset$, то $P_-=0$ и потому 


$B_k=0$ $\forall k\le 0$. 


Если $\sigma_-\ne 0$, то $P_-\ne 0$ и в этом 


случае рассмотрим число $r_-\in(0,1)$, удовлетворяющее условию 


$(1-r_-)\|\cal D^{-1}\|<1$. Пусть 


$\Bbb T(r_-)=\{\lambda\in\Bbb C:|\lambda|=r_-\}$. 


Включение 


$\Bbb T(r_-)\subset\rho(A)$ следует из представления 


$\theta r_-I-A=(\theta I-A)[I-\theta(1-r_-)R(\theta,A)]$, 


$\theta\in\Bbb T$, 


и равенств $\|R(\theta,A)\|=\|\cal D^{-1}\|$, $\theta\in\Bbb T$, вытекающих 


из (9). При этом 


имеет место оценка 


$\|R(\lambda,A)\|\le\|\cal D^{-1}\|(1-\delta_-\|\cal D^{-1}\|)^{-1}$, 


где $\delta_-=1-r_-$ и $\lambda\in\Bbb T(r_-)$. Отсюда и из (16) следует 


представление 


$$ 


B_k=\frac{1}{2\pi i}\int_{\Bbb T(r_-)}\lambda^{-k} 


R(\lambda,A)d\lambda,\quad k\le -1, 


$$ 


из которого вытекают оценки (13). 


 


Аналогично устанавливаются оценки (14). В этом случае рассматривается 


число $r_+>1$, удовлетворяющее условию 


$\delta_+\|\cal D^{-1}\|<1$, $\delta_+=r_+-1$, и используется окружность 


$\Bbb T(r_+)$ 


для интегрального представления операторов $B_k$, $k\ge 1$. 


\end{pf} 


 


В условиях следующей теоремы рассматривается последовательность функций 


$B_k(g)=B(g)B(g-g_0)\dots B(g-(k-1)g_0)$, $k\ge 1$, 


из пространства $C_S(G,\End X)$. 


 


\begin{thm} Для того чтобы оператор $\cal D=I-A$ был обратим, необходимо и 


достаточно, 


чтобы существовала проекторозначная функция $P:C_S(G,\End X)$ со следующими 


свойствами 


{\em (}в этом случае будем говорить, что последовательность функций обладает 


дискретной дихотомией; см. {\em (\cite{Hen},} определение {\em 7.6.4)):} 


\begin{itemize} 


\item[1)] $B(g)P(g-g_0)=P(g)B(g)$ $\forall g\in G;$ 


\item[2)] найдется число $k\in\Bbb N$ такое, что 


$\sup\limits_{g\in G}\|P(g)B_k(g)\|<1;$ 


\item[3)] существуют числа $m\in\Bbb N$ и $q>1$, для которых 


$$ 


\|Q(g)B_m(g)x\|=\|B_m(g)Q(g-(m-1)g_0)x\| \ge 


q\|Q(g-(m-1)g_0)x\| 


$$ 


при всех $x\in X$, $g\in G$ и $B_m(g)$ осуществляет изоморфизм образа 


$\im Q(g-(m-1)g_0)$ 


проектора $Q(g-(m-1)g_0)$ и $\im Q(g)$ при любом $g\in G$ 


$(Q(g)=I-P(g)$, $g\in G)$. 


\end{itemize} 


\end{thm} 


 


\begin{pf} Необходимость условий непосредственно следует из теоремы 2. 


Существование 


функции $P$ с указанными свойствами означает, что проектор 


$(P_-x)(g)=P(g)x(g)$, $g\in G$, $x\in\cal F$,


перестановочен с оператором $A$ (это следует из условия 1)), спектр сужения 


оператора $A$ на подпространство $\cal F_-=\im P_-$ лежит в круге 


$\{\lambda\in\Bbb C:|\lambda|<1\}$ (что 


вытекает из 


условия 2)), оператор 


$A_+=A\mid\im P_+$, $P_+=I-P_-$, обратим и 


$\sigma(A_+^{-1})\subset\{\lambda\in\Bbb C:|\lambda|<1\}$. Поэтому 


$\sigma(A)\cap\Bbb T=\emptyset$, т.\,е. 


обратим оператор $\cal D$. 


\end{pf} 


 


\begin{cornonum} Если оператор 


$\cal D=I-A\in\End\cal F$ обратим, то его обратный определяется 


формулой 


$$ 


(\cal D^{-1}x)(g)=\sum_{k\in\Bbb Z}\cal D_k(g)x(g-kg_0), 


\quad g\in G,\quad x\in\cal F, 


$$ 


где 


$\cal D_k(g)=B_k(g)P(g-kg_0)=P(g)B_k(g)$, $g\in G$, $k\ge 1$, 


$\cal D_0(g)=P(g)$, $\cal D_k(g)x=0$ $\forall x\in\im Q(g-(k-1)g_0)$ 


и $\cal D_k(g)=(B_{-k}(g)\mid\im Q(g-kg_0))^{-1}$ 


на $\im Q(g-kg_0)$ при $k\le -1$. 


\end{cornonum} 


 


Полученные в теоремах 1--3 результаты применим к исследованию линейного 


дифференциального оператора 


$$ 


\cal L=-\frac{d}{dt}+A(t):D(\cal L)\subset\cal F\to\cal F= 


\cal F(\Bbb R,X) 


$$ 


в предположении, что семейство замкнутых линейных операторов 


$A(t):D(A(t))\subset X\to X$, $t\in\Bbb R$, порождает 


корректную задачу Коши для дифференциального уравнения 


$\dfrac{dx}{dt}=A(t)x$ (см. \cite{Kre}). 


Пусть $\{\cal U(t,s),\;s\le t\}$ --- семейство эволюционных операторов для 


этого уравнения. 


 


Будем считать функцию $x\in\cal F$ принадлежащей области определения 


$D(\cal L)$ 


оператора $\cal L$, 


если существует функция $f\in\cal F$ такая, что для почти всех $s\le t$ 


имеют место равенства 


$$ 


x(t)=\cal U(t,s)x(s)-\int^t_s \cal U(t,\tau)f(\tau)d\tau. 


$$ 


В этом случае полагаем $\cal L x=f$. При исследовании оператора $\cal L$ 


возникают проблемы, 


связанные, в первую очередь, с его неограниченностью. Чтобы их снять,


семейству $\{\cal U(t,s),\;s\le t\}$ сопоставим полугруппу разностных 


операторов $\{T(t),\;t\ge 0\}$ из банаховой алгебры $\End\cal F$, имеющих 


вид 


$$ 


(T(t)x)(s)=\cal U(s,s-t)x(s-t),\quad 


s\in\Bbb R,\quad x\in\cal F,\quad t\ge 0. 


$$ 


Отметим два следующих утверждения из статьи \cite{Bask2}. 


 


\begin{thm} Оператор $\cal L$ является производящим оператором полугруппы 


операторов $\{T(t)$,


%//     join ^^^


$t\ge 0\}$ в любом из банаховых пространств 


$L_p$, $p\in[1,\infty)$, $C_0$. 


\end{thm} 


 


\begin{thm} Спектр $\sigma(\cal L)$ оператора 


$\cal L:D(\cal L)\subset\cal F\to\cal F$ не зависит от выбора 


пространства $\cal F$. 


\end{thm} 


 


Имеет место следующая теорема об отображении спектра, которая 


неверна для произвольной сильно непрерывной полугруппы линейных операторов 


(см. \cite{Hil}, гл.\,16). 


 


\begin{thm} Для любого числа $t>0$ справедливо равенство 


$$ 


\sigma(T(t))\setminus\{0\}=\exp \sigma(\cal L)t= 


\exp\{ \lambda t:\lambda\in\sigma(\cal L)\}. 


$$ 


\end{thm} 


 


\begin{pf} Поскольку оператор $\cal L-\lambda I$, $\lambda\in\Bbb C$, 


является производящим оператором полугруппы 


операторов 


$\{T(t)\exp\lambda t,\;t\ge 0\}$, то достаточно доказать, что оператор 


$\cal L$ 


обратим тогда и только тогда, когда $1\notin\sigma(T(t))$. 


 


Если оператор $\cal L$ обратим, то в силу теоремы 5 он обратим в банаховом 


пространстве $C$. В этом случае семейство эволюционных операторов 


$\{\cal U(s,\tau),\;\tau\le s\}$ 


допускает 


экспоненциальную дихотомию (см. \cite{Lev}, гл.\,X). Следовательно, 


семейство функций 


$B_k(s)=\cal U(s,s-kt)$, $k\ge 1$, допускает дискретную экспоненциальную 


дихотомию (см. 


формулировку теоремы 3). Поэтому из теоремы 3 следует обратимость оператора 


$I-T(t)$, т.\,е. $1\notin\sigma(T(t))$. Из следствия теоремы 3 получаем 


формулу 


$$ 


((T(t)-I)^{-1}x)(s)=\sum_{k\in\Bbb Z} 


G(s,s+kt)x(s+kt),\quad x\in\cal F, 


$$ 


где 


$G(s,\tau)=\{\cal U(s,\tau)p(\tau),\; \tau\le s;\;


\cal U(s,\tau)Q(\tau),\; \tau>s\}$, $Q(\tau)=I-P(\tau)$ 


и проекторозначная функция $P:\Bbb R\to\End X$ взята из определения 


экспоненциальной дихотомии семейства эволюционных операторов (см. 


\cite{Lev}, гл.\,X). 


 


Обратно, если $1\notin\sigma(T(t))$, то из (\cite{Hil}, гл.\,XVI) следует, 


что $0\notin\sigma(\cal L)$, если полугруппа $\{T(t),\;t\ge 0\}$


сильно непрерывна (т.\,е. если 


$\cal F=L_p$, $p\in[1,\infty)$, или $\cal F=C_0$). Осталось еще раз 


воспользоваться теоремой 5. 


\end{pf} 


 


Доказанная теорема позволяет свести изучение ряда спектральных свойств 


оператора $\cal L$ к изучению соответствующих свойств разностных 


операторов. 


 


\begin{thebibliography}{99} 


 


\bibitem{Lev} 


Левитан Б.М., Жиков В.В. {\em Почти периодические функции и дифференциальные 


уравнения}. -- М.: Изд-во МГУ, 1978. -- 204 с. 


 


\bibitem{Hew} 


Хьюитт Э., Росс К. {\em Абстрактный гармонический анализ}. -- М.: Наука, Мир, 


1975. -- T.\,1,~2. -- 657 с., 901 с. 


 


\bibitem{Dun} 


Данфорд Н., Шварц Дж.Т. {\em Линейные операторы}. -- М.: Ин. лит.,


1962. -- Т.\,1. -- 895 с. 


 


\bibitem{Bou} 


Бурбаки Н. {\em Спектральная теория}. -- М.: Мир, 1972. -- 183 с. 


 


\bibitem{Kur} 


Курбатов В.Г. {\em Линейные дифференциально-разностные уравнения}. -- Воронеж: 


Изд-во ВГУ, 1990. -- 164 с. 


 


\bibitem{Muh} 


Мухамадиев Э. {\em Об обратимости функциональных операторов в пространстве 


ограниченных на оси функций} // Матем. заметки. -- 1972. -- Т.11. 


-- С.269--274. 


 


\bibitem{Slu} 


Слюсчарчук В.Е. {\em Обратимость неавтономных дифференциально-функциональных 


операторов} // Матем. сб. -- 1986. -- Т.130. -- \No\,1. -- С.86--104. 


 


\bibitem{Bask1} 


Баскаков А.Г. {\em Абстрактный гармонический анализ и асимптотические


оценки элементов обратных матриц} // Матем. заметки. -- 1992. -- Т.52.


-- \No\,2. 


-- С.17--26. 


 


\bibitem{Ant} 


Антоневич А.Б. {\em Линейные функциональные уравнения: операторный подход}. -- 


Минск: 1988. -- 232 с. 


 


\bibitem{Hil} 


Хилле Э., Филлипс Р. {\em Функциональный анализ и полугруппы}. -- М.: Ин. лит., 


1962. -- 829 с. 


 


\bibitem{Zab} 


Забрейко П.П., Нгуен Вань Мин. {\em Экспоненциальная дихотомия и интегральные 


многообразия в теории потоков и их применения} // Докл. РАН. -- 1992. -- 


Т.324. -- \No\,3. -- С.515-518. 


 


\bibitem{Kre} 


Крейн С.Г. {\em Линейные дифференциальные уравнения в банаховом пространстве}.


--  М.: Наука, 1967. -- 464 с. 


 


\bibitem{Bask2} 


Баскаков А.Г. {\em Спектральный анализ линейных дифференциальных операторов и 


полугруппы разностных операторов} // Докл. РАН. -- 1995. -- Т.343. -- 


\No\,3. -- С.295--298. 


 


\bibitem{Hen} 


Хенри Д. {\em Геометрическая теория полулинейных параболических уравнений}. -- 


М.: Мир, 1985. -- 376 с. 


 


 


\end{thebibliography} 


 


\vskip 2cm 


 


\post{\it Воронежский государственный университет}{\it Поступила} 


\post{}{07.02.1995} 


 


\end{document} 


